% ======================================================================
% Chapter 5: Elementary Processes of Quantum Electrodynamics
% (book pages 131-173; PDF pages 154-197)
% ======================================================================
\chapter{Elementary Processes of Quantum Electrodynamics}

Finally, after three long chapters of formalism, we are ready to perform some
real relativistic calculations, to begin working out the predictions of Quantum
Electrodynamics.  First we will return to the process considered in Chapter~1,
the annihilation of an electron-positron pair into a pair of heavier fermions.
We will study this paradigm process in extreme detail in the next three
sections, then do a few more simple QED calculations in Sections~5.4 and~5.5.
The problems at the end of the chapter treat several additional QED processes.
More complete surveys of QED can be found in the books of Jauch and Rohrlich
(1976) and of Berestetskii, Lifshitz, and Pitaevskii (1982).

\section{$e^{+}e^{-}\to\mu^{+}\mu^{-}$: Introduction}

The reaction $e^{+}e^{-}\to\mu^{+}\mu^{-}$ is the simplest of all QED processes,
but also one of the most important in high-energy physics.  It is fundamental
to the understanding of all reactions in $e^{+}e^{-}$ colliders, and is in fact
used to calibrate such machines.  The related process $e^{+}e^{-}\to q\bar{q}$
(a quark-antiquark pair) is extraordinarily useful in determining the
properties of elementary particles.

In this section we will compute the unpolarized cross section for
$e^{+}e^{-}\to\mu^{+}\mu^{-}$ to lowest order.  In Chapter~1 we used elementary
arguments to guess the answer (Eq.~\eqref{eq:1.8}) in the limit where all the
fermions are massless.  We now relax that restriction and retain the muon
mass in the calculation.  Retaining the electron mass as well would be easy
but pointless, since the ratio $m_e/m_\mu \approx 1/200$ is much smaller than
the fractional error introduced by neglecting higher-order terms in the
perturbation series.

Using the Feynman rules from Section~4.8, we can at once draw the diagram and
write down the amplitude for our process:
\begin{equation}
  i\mathcal{M} = \bar{v}^{s'}(p')(-ie\gamma^\mu)u^s(p)\
    \frac{-ig_{\mu\nu}}{q^2}\ \bar{u}^{r}(k)(-ie\gamma^\nu)v^{r'}(k').
\end{equation}
Rearranging this slightly and leaving the spin superscripts implicit, we have
\begin{equation}
  i\mathcal{M}\bigl( e^{-}(p)e^{+}(p') \to \mu^{-}(k)\mu^{+}(k') \bigr)
    = \frac{(-ie)^2}{q^2}\ \bigl( \bar{v}(p')\gamma^\mu u(p) \bigr)\
      \bigl( \bar{u}(k)\gamma_\mu v(k') \bigr) .
  \label{eq:5.1}
\end{equation}
This answer for the amplitude $\mathcal{M}$ is simple, but not yet very
illuminating.  To compute the differential cross section, we need an
expression for $|\mathcal{M}|^2$, so we must find the complex conjugate of
$\mathcal{M}$.  A bi-spinor product such as $\bar{v}\gamma^\mu u$ can be
complex-conjugated as follows:
\begin{equation}
  \bigl( \bar{v}\gamma^\mu u \bigr)^{*}
    = u^{\dagger}\gamma^{\mu\dagger}(\gamma^0)^{\dagger}v
    = \bar{u}\gamma^\mu v .
\end{equation}
(This is another advantage of the ``bar'' notation.)  Thus the squared matrix
element is
\begin{equation}
  |\mathcal{M}|^2 = \frac{e^4}{q^4}\ \bigl( \bar{v}(p')\gamma^\mu u(p)
    \bar{u}(p)\gamma^\nu v(p') \bigr)\
    \bigl( \bar{u}(k)\gamma_\mu v(k')\ \bar{v}(k')\gamma_\nu u(k) \bigr).
  \label{eq:5.2}
\end{equation}
At this point we are still free to specify any particular spinors $u^s(p)$,
$v^{s'}(p')$, and so on, corresponding to any desired spin states of the
fermions.  In actual experiments, however, it is difficult (though not
impossible) to retain control over spin states; one would have to prepare the
initial state from polarized materials and/or analyze the final state using
spin-dependent multiple scattering.  In most experiments the electron and
positron beams are unpolarized, so the measured cross section is an average
over the electron and positron spins $s$ and $s'$.  Muon detectors are
normally blind to polarization, so the measured cross section is a sum over
the muon spins $r$ and $r'$.

The expression for $|\mathcal{M}|^2$ simplifies considerably when we throw
away the spin information.  We want to compute
\begin{equation}
  \frac{1}{4}\sum_{s,s'}\sum_{r,r'}|\mathcal{M}|^2 .
\end{equation}
The spin sums can be performed using the completeness relations from
Section~3.3:
\begin{equation}
  \sum_s u^s(p)\bar{u}^s(p) = \slashed{p} + m;\qquad
  \sum_s v^s(p)\bar{v}^s(p) = \slashed{p} - m .
  \label{eq:5.3}
\end{equation}
Working with the first half of \eqref{eq:5.2}, and writing in spinor indices
so we can freely move the $v$ next to the $\bar{v}$, we have
\begin{equation}
  \begin{aligned}
    \sum_{s,s'} (\slashed{p}' - m)_{da}\gamma^\mu_{ab}(\slashed{p} + m)_{bc}
      \gamma^\nu_{cd} &= \mathrm{tr}\bigl[
        (\slashed{p}' - m)\gamma^\mu(\slashed{p} + m)\gamma^\nu \bigr].
  \end{aligned}
\end{equation}
Evaluating the second half of \eqref{eq:5.2} in the same way, we arrive at
the desired simplification:
\begin{equation}
  \frac{1}{4}\sum_{\text{spins}}|\mathcal{M}|^2
    = \frac{e^4}{4q^4}\ \mathrm{tr}\bigl[
      (\slashed{p}' - m_e)\gamma^\mu(\slashed{p} + m_e)\gamma^\nu \bigr]\
      \mathrm{tr}\bigl[ (\slashed{k} + m_\mu)\gamma_\mu(\slashed{k}' - m_\mu)
        \gamma_\nu \bigr] .
  \label{eq:5.4}
\end{equation}
The spinors $u$ and $v$ have disappeared, leaving us with a much cleaner
expression in terms of $\gamma$ matrices.  This trick is very general: Any QED
amplitude involving external fermions, when squared and summed or averaged
over spins, can be converted in this way to traces of products of Dirac
matrices.

\subsection{Trace Technology}

This last step would hardly be an improvement if the traces had to be
laboriously computed by brute force.  But Feynman found that they could be
worked out easily by appealing to the algebraic properties of the $\gamma$
matrices.  Since the evaluation of such traces occurs so often in QED
calculations, it is worthwhile to pause and attack the problem
systematically, once and for all.

We would like to evaluate traces of products of $n$ gamma matrices, where
$n = 0, 1, 2, \ldots$.  (For the present problem we need $n = 2, 3, 4$.)  The
$n=0$ case is fairly easy: $\mathrm{tr}\,1 = 4$.  The trace of one $\gamma$
matrix is also easy.  From the explicit form of the matrices in the chiral
representation, we have
\begin{equation}
  \mathrm{tr}\,\gamma^\mu = 0 .
\end{equation}
It is useful to prove this result in a more abstract way, which generalizes to
an arbitrary odd number of $\gamma$ matrices:
\begin{equation}
  \begin{aligned}
    \mathrm{tr}\,\gamma^\mu &= \mathrm{tr}\,\gamma^5\gamma^5\gamma^\mu
      &&\text{since } (\gamma^5)^2 = 1\\
    &= -\mathrm{tr}\,\gamma^5\gamma^\mu\gamma^5 &&\text{since }
      \{\gamma^\mu, \gamma^5\} = 0\\
    &= -\mathrm{tr}\,\gamma^5\gamma^5\gamma^\mu
      &&\text{using cyclic property of trace}\\
    &= -\mathrm{tr}\,\gamma^\mu .
  \end{aligned}
\end{equation}
Since the trace of $\gamma^\mu$ is equal to minus itself, it must vanish.  For
$n$ $\gamma$-matrices we would get $n$ minus signs in the second step (as we
move the second $\gamma^5$ all the way to the right), so the trace must vanish
if $n$ is odd.

To evaluate the trace of two $\gamma$ matrices, we again use the
anticommutation properties and the cyclic property of the trace:
\begin{equation}
  \begin{aligned}
    \mathrm{tr}\,\gamma^\mu\gamma^\nu
      &= \mathrm{tr}(2g^{\mu\nu}\cdot 1 - \gamma^\nu\gamma^\mu)
      &&\text{(anticommutation)}\\
    &= 8g^{\mu\nu} - \mathrm{tr}\,\gamma^\mu\gamma^\nu
      &&\text{(cyclicity)}
  \end{aligned}
\end{equation}
Thus $\mathrm{tr}\,\gamma^\mu\gamma^\nu = 4g^{\mu\nu}$.  The trace of any even
number of $\gamma$ matrices can be evaluated in the same way: Anticommute the
first $\gamma$ matrix all the way to the right, then cycle it back to the
left.  Thus for the trace of four $\gamma$ matrices, we have
\begin{equation}
  \begin{aligned}
    \mathrm{tr}\,(\gamma^\mu\gamma^\nu\gamma^\rho\gamma^\sigma)
    &= \mathrm{tr}(2g^{\mu\nu}\gamma^\rho\gamma^\sigma
      - \gamma^\nu\gamma^\mu\gamma^\rho\gamma^\sigma)\\
    &= \mathrm{tr}(2g^{\mu\nu}\gamma^\rho\gamma^\sigma
      - 2g^{\mu\rho}\gamma^\nu\gamma^\sigma
      + 2g^{\mu\sigma}\gamma^\nu\gamma^\rho
      - \gamma^\nu\gamma^\rho\gamma^\sigma\gamma^\mu) .
  \end{aligned}
\end{equation}
Using the cyclic property on the last term and bringing it to the left-hand
side, we find
\begin{equation}
  \mathrm{tr}(\gamma^\mu\gamma^\nu\gamma^\rho\gamma^\sigma)
    = g^{\mu\nu}\mathrm{tr}\,\gamma^\rho\gamma^\sigma
      - g^{\mu\rho}\mathrm{tr}\,\gamma^\nu\gamma^\sigma
      + g^{\mu\sigma}\mathrm{tr}\,\gamma^\nu\gamma^\rho
    = 4\bigl( g^{\mu\nu}g^{\rho\sigma} - g^{\mu\rho}g^{\nu\sigma}
      + g^{\mu\sigma}g^{\nu\rho} \bigr) .
\end{equation}
In this manner one can always reduce a trace of $n$ $\gamma$-matrices to a sum
of traces of $(n-2)$ $\gamma$-matrices.  The case $n=6$ is easy to work out,
but has fifteen terms (the number of ways of grouping the six indices in pairs
to make terms of the form $g^{\mu\nu}g^{\rho\sigma}g^{\alpha\beta}$).
Fortunately, we will not need it in this book.  (If you ever do need to
evaluate such complicated traces, it may be easier to learn to use one of the
several computer programs that can perform symbolic manipulations on Dirac
matrices.)

Starting in Section~5.2, we will often need to evaluate traces involving
$\gamma^5$.  Since $\gamma^5 = i\gamma^0\gamma^1\gamma^2\gamma^3$, the trace
of $\gamma^5$ times any odd number of other $\gamma$ matrices is zero.  It is
also easy to show that the trace of $\gamma^5$ itself is zero:
\begin{equation}
  \mathrm{tr}\,\gamma^5 = \mathrm{tr}(\gamma^0\gamma^0\gamma^5)
    = -\mathrm{tr}(\gamma^0\gamma^5\gamma^0)
    = -\mathrm{tr}(\gamma^0\gamma^0\gamma^5) = -\mathrm{tr}\,\gamma^5 .
\end{equation}
The same trick works for $\mathrm{tr}(\gamma^\mu\gamma^\nu\gamma^5)$, if we
insert two factors of $\gamma^\alpha$ for some $\alpha$ different from both
$\mu$ and $\nu$.  The first nonvanishing trace involving $\gamma^5$ contains
four other $\gamma$ matrices.  In this case the trick still works unless every
$\gamma$ matrix appears, so
\begin{equation}
  \mathrm{tr}(\gamma^\mu\gamma^\nu\gamma^\rho\gamma^\sigma\gamma^5) = 0
  \qquad\text{unless } (\mu\nu\rho\sigma) \text{ is some permutation of }
    (0123).
\end{equation}
From the anticommutation rules it also follows that interchanging any two of
the indices simply changes the sign of the trace, so
$\mathrm{tr}(\gamma^\mu\gamma^\nu\gamma^\rho\gamma^\sigma\gamma^5)$ must be
proportional to $\epsilon^{\mu\nu\rho\sigma}$.  The overall constant turns out
to be $-4i$, as you can easily check by plugging in
$(\mu\nu\rho\sigma) = (0123)$.

Here is a summary of the trace theorems, for convenient reference:
\begin{equation}
  \begin{aligned}
    \mathrm{tr}(1) &= 4\\
    \mathrm{tr}(\text{any odd \# of }\gamma\text{'s}) &= 0\\
    \mathrm{tr}(\gamma^\mu\gamma^\nu) &= 4g^{\mu\nu}\\
    \mathrm{tr}(\gamma^\mu\gamma^\nu\gamma^\rho\gamma^\sigma)
      &= 4\bigl( g^{\mu\nu}g^{\rho\sigma} - g^{\mu\rho}g^{\nu\sigma}
        + g^{\mu\sigma}g^{\nu\rho} \bigr)\\
    \mathrm{tr}(\gamma^5) &= 0\\
    \mathrm{tr}(\gamma^\mu\gamma^\nu\gamma^5) &= 0\\
    \mathrm{tr}(\gamma^\mu\gamma^\nu\gamma^\rho\gamma^\sigma\gamma^5)
      &= -4i\epsilon^{\mu\nu\rho\sigma}
  \end{aligned}
  \label{eq:5.5}
\end{equation}
Expressions resulting from use of the last formula can be simplified by means
of the identities
\begin{equation}
  \epsilon_{\alpha\beta\gamma\delta}\epsilon^{\alpha\beta\gamma\delta}
    = -24;\qquad
  \epsilon_{\alpha\beta\mu\nu}\epsilon^{\alpha\beta\rho\sigma}
    = -2\bigl( \delta^\rho_\mu\delta^\sigma_\nu
      - \delta^\rho_\nu\delta^\sigma_\mu \bigr) .
  \label{eq:5.6}
\end{equation}
All of these can be derived by first appealing to symmetry arguments, then
evaluating one special case to determine the overall constant.

Another useful identity allows one to reverse the order of all the $\gamma$
matrices inside a trace:
\begin{equation}
  \mathrm{tr}(\gamma^\mu\gamma^\nu\gamma^\rho\gamma^\sigma)
    = \mathrm{tr}(\gamma^\sigma\gamma^\rho\gamma^\nu\gamma^\mu) .
  \label{eq:5.7}
\end{equation}

We can now complete the evaluation of \eqref{eq:5.4}.  Dropping the electron
mass (setting $m_e = 0$), we find
\begin{equation}
  \frac{1}{4}\sum_{\text{spins}}|\mathcal{M}|^2
    = \frac{e^4}{4q^4}\ \bigl[ 4\bigl( p'^{\mu}p^{\nu} - p'\cdot p\,g^{\mu\nu}
      + p'^{\nu}p^{\mu} \bigr) \bigr]\ \bigl[ 4\bigl( k_\mu k'_\nu
        - k\cdot k'\,g_{\mu\nu} + k_\nu k'_\mu \bigr) \bigr].
  \label{eq:5.8}
\end{equation}
Contracting the indices and using $q^2 = (p+p')^2 = 2p\cdot p'$, we obtain
\begin{equation}
  \frac{1}{4}\sum_{\text{spins}}|\mathcal{M}|^2
    = \frac{8e^4}{q^4}\ \bigl[ (p\cdot k)(p'\cdot k') + (p\cdot k')(p'\cdot k)
      + m_\mu^2\,p\cdot p' \bigr] .
  \label{eq:5.9}
\end{equation}
Specializing to the center-of-mass frame with $E$ the energy of the electron
(and positron) and $E_\mu$ the energy of the muon (and antimuon), we have
$p = (E, \mathbf{p})$, $p' = (E, -\mathbf{p})$, $k = (E_\mu, \mathbf{k})$,
$k' = (E_\mu, -\mathbf{k})$, and $q^2 = 4E^2$.  A straightforward computation
gives
\begin{equation}
  \frac{1}{4}\sum_{\text{spins}}|\mathcal{M}|^2
    = \frac{e^4}{4E^4}\ \bigl[ (E + E_\mu)^2 + (E\cos\theta + E_\mu)^2
      + m_\mu^2 \bigr] .
  \label{eq:5.10}
\end{equation}
The cross section is then obtained from the general formula
\eqref{eq:4.82}:
\begin{equation}
  \Bigl( \frac{d\sigma}{d\Omega} \Bigr)_{\mathrm{CM}}
    = \frac{1}{64\pi^2 E_{\mathrm{cm}}^2}\ \frac{1}{4}
      \sum_{\text{spins}}|\mathcal{M}|^2 .
  \label{eq:5.11}
\end{equation}
Substituting \eqref{eq:5.10} and simplifying with
$E_{\mathrm{cm}} = 2E$, $E_\mu = E_{\mathrm{cm}}/2$, we find
\begin{equation}
  \Bigl( \frac{d\sigma}{d\Omega} \Bigr)_{\mathrm{CM}}
    = \frac{\alpha^2}{4E_{\mathrm{cm}}^2}\ \sqrt{1 - \frac{m_\mu^2}{E^2}}
      \Bigl[ 1 + \cos^2\theta + \frac{m_\mu^2}{E^2}\sin^2\theta \Bigr] ,
  \label{eq:5.12}
\end{equation}
where $\alpha = e^2/4\pi$.  Integrating over $\theta$ and $\phi$ gives the
total cross section,
\begin{equation}
  \sigma_{\mathrm{total}}
    = \frac{4\pi\alpha^2}{3E_{\mathrm{cm}}^2}\ \sqrt{1 - \frac{m_\mu^2}{E^2}}
      \Bigl[ 1 + \frac{m_\mu^2}{2E^2} \Bigr] .
  \label{eq:5.13}
\end{equation}
In the high-energy limit where $E \gg m_\mu$, these formulae reduce to those
given in Chapter~1:
\begin{equation}
  \frac{d\sigma}{d\Omega} \to \frac{\alpha^2}{4E_{\mathrm{cm}}^2}
    (1 + \cos^2\theta);\qquad
  \sigma_{\mathrm{total}} \to \frac{4\pi\alpha^2}{3E_{\mathrm{cm}}^2} .
  \label{eq:5.14}
\end{equation}
Note that these expressions have the correct dimensions of cross sections.
In the high-energy limit, $E_{\mathrm{cm}}$ is the only dimensionful quantity
in the problem, so dimensional analysis dictates that
$\sigma_{\mathrm{total}} \sim E_{\mathrm{cm}}^{-2}$.  Since we knew from the
beginning that $\sigma_{\mathrm{total}} \propto \alpha^2$, we only had to work
to get the factor of $4\pi/3$.

The energy dependence of the total cross-section formula \eqref{eq:5.13}
near threshold is shown in Fig.~\ref{fig:5.1}.  Of course the cross section is
zero for $E_{\mathrm{cm}} < 2m_\mu$.  It is interesting to compare the shape of
the actual curve to the shape one would obtain if $|\mathcal{M}|^2$ did not
depend on energy, that is, if all the energy dependence came from the
phase-space factor $|\mathbf{k}|/E$.  To test Quantum Electrodynamics, an
experiment must be able to resolve deviations from the naive phase-space
prediction.  Experimental results from pair production of both $\mu$ and
$\tau$ leptons confirm that these particles behave as QED predicts.
Figure~\ref{fig:5.2} compares formula \eqref{eq:5.13} to experimental
measurements of the $\tau^{+}\tau^{-}$ threshold.

\begin{figure}[t]
  \centering
  \includegraphics[width=0.6\textwidth]{fig5.1.png}
  \caption{Energy dependence of the total cross section for
    $e^{+}e^{-}\to\mu^{+}\mu^{-}$ compared to ``phase space'' energy
    dependence.}
  \label{fig:5.1}
\end{figure}

\begin{figure}[t]
  \centering
  \includegraphics[width=0.7\textwidth]{fig5.2.png}
  \caption{The ratio $\sigma(e^{+}e^{-}\to\tau^{+}\tau^{-})/\sigma(e^{+}e^{-}
    \to\mu^{+}\mu^{-})$ of measured cross sections near the threshold for
    $\tau^{+}\tau^{-}$ pair-production, as measured by the DELCO
    collaboration, W.~Bacino, et.\ al., \emph{Phys.\ Rev.\ Lett.} \textbf{41},
    13 (1978).  Only a fraction of $\tau$ decays are included, hence the small
    overall scale.  The curve shows a fit to the theoretical formula
    \eqref{eq:5.13}, with a small energy-independent background added.  The
    fit yields $m_\tau = 1782\frac{+2}{-2}$ MeV.}
  \label{fig:5.2}
\end{figure}

Before discussing our result further, let us pause to summarize how we
obtained it.  The method extends in a straightforward way to the calculation
of unpolarized cross sections for other QED processes.  The general procedure
is as follows:
\begin{enumerate}
  \item Draw the diagram(s) for the desired process.
  \item Use the Feynman rules to write down the amplitude $\mathcal{M}$.
  \item Square the amplitude and average or sum over spins, using the
    completeness relations \eqref{eq:5.3}.  (For processes involving photons
    in the final state there is an analogous completeness relation, derived
    in Section~5.5.)
  \item Evaluate traces using the trace theorems \eqref{eq:5.5}; collect terms
    and simplify the answer as much as possible.
  \item Specialize to a particular frame of reference, and draw a picture of
    the kinematic variables in that frame.  Express all 4-momentum vectors in
    terms of a suitably chosen set of variables such as $E$ and $\theta$.
  \item Plug the resulting expression for $|\mathcal{M}|^2$ into the
    cross-section formula \eqref{eq:4.79}, and integrate over phase-space
    variables that are not measured to obtain a differential cross section in
    the desired form.  (In our case these integrations were over the
    constrained momenta $\mathbf{k}'$ and $|\mathbf{k}|$, and were performed
    in the derivation of Eq.~(4.84).)
\end{enumerate}
While other calculations (especially those involving loop diagrams) often
require additional tricks, nearly every QED calculation will involve the basic
procedures outlined here.

\subsection{Production of Quark-Antiquark Pairs}

The asymptotic energy dependence of the $e^{+}e^{-}\to\mu^{+}\mu^{-}$
cross-section formula sets the scale for all $e^{+}e^{-}$ annihilation cross
sections.  A particularly important example is the cross section for
\begin{equation}
  e^{+}e^{-} \to \text{hadrons},
\end{equation}
that is, the total cross section for production of any number of strongly
interacting particles.

In our current understanding of the strong interactions, given by the theory
called Quantum Chromodynamics (QCD), all hadrons are composed of Dirac
fermions called quarks.  Quarks appear in a variety of types, called flavors,
each with its own mass and electric charge.  A quark also carries an
additional quantum number, color, which takes one of three values.  Color
serves as the ``charge'' of QCD, as we will discuss in Chapter~17.

According to QCD, the simplest $e^{+}e^{-}$ process that ends in hadrons is
\begin{equation}
  e^{+}e^{-} \to q\bar{q},
\end{equation}
the annihilation of an electron and a positron, through a virtual photon, into
a quark-antiquark pair.  After they are created, the quarks interact with one
another through their strong forces, producing more quark pairs.  Eventually
the quarks and antiquarks combine to form some number of mesons and baryons.

To adapt our results for muon production to handle the case of quarks, we must
make three modifications:
\begin{enumerate}
  \item Replace the muon charge $e$ with the quark charge $Q|e|$.
  \item Count each quark three times, one for each color.
  \item Include the effects of the strong interactions of the produced quark
    and antiquark.
\end{enumerate}
The first two changes are easy to make.  For the first, it is simply necessary
to know the masses and charges of each flavor of quark.  For $u$, $c$, and $t$
quarks we have $Q = 2/3$, while for $d$, $s$, and $b$ quarks we have
$Q = -1/3$.  The cross-section formulae are proportional to the square of the
charge of the final-state particle, so we can simply insert a factor of $Q^2$
into any of these formulae to obtain the cross section for production of any
particular variety of quark.  Counting colors is necessary because experiments
measure only the total cross section for production of all three colors.  (The
hadrons that are actually detected are colorless.)  In any case, this counting
is easy: Just multiply the answer by 3.

If you know a little about the strong interaction, however, you might think
this is all a big joke.  Surely the third modification is extremely difficult
to make, and will drastically alter the predictions of QED.  The amazing truth
is that in the high-energy limit, the effect of the strong interaction on the
quark production process can be completely neglected.  As we will discuss in
Part~III, the only effect of the strong interaction (in this limit) is to
dress up the final-state quarks into bunches of hadrons.  This simplification
is due to a phenomenon called asymptotic freedom; it played a crucial role in
the identification of Quantum Chromodynamics as the correct theory of the
strong force.

Thus in the high-energy limit, we expect the cross section for the reaction
$e^{+}e^{-}\to q\bar{q}$ to approach $3\cdot Q^2\cdot 4\pi\alpha^2/3E^2_{\mathrm{cm}}$.
It is conventional to define
\begin{equation}
  R \equiv \frac{\sigma(e^{+}e^{-}\to\text{hadrons})}
    {\sigma(e^{+}e^{-}\to\mu^{+}\mu^{-})}
  = \frac{4\pi\alpha^2}{3E_{\mathrm{cm}}^2}
    \quad\text{in units of } R
    = \Bigl( \frac{E_{\mathrm{cm}}}{1\ \mathrm{GeV}} \Bigr)^2 \cdot
      \frac{86.8\ \mathrm{nbarns}}{4\pi\alpha^2/3E_{\mathrm{cm}}^2}
      \cdot\ \cdots .
  \label{eq:5.15}
\end{equation}
The value of a cross section in units of $R$ is therefore its ratio to the
asymptotic value of the $e^{+}e^{-}\to\mu^{+}\mu^{-}$ cross section predicted
by Eq.~\eqref{eq:5.14}.  Experimentally, the easiest quantity to measure is
the total rate for production of all hadrons.  Asymptotically, we expect
\begin{equation}
  \frac{\sigma(e^{+}e^{-}\to\text{hadrons})}{R}
    \to 3\cdot\Bigl( \sum_f Q_f^2 \Bigr),
  \qquad E_{\mathrm{cm}} \to \infty ,
  \label{eq:5.16}
\end{equation}
where the sum runs over all quarks whose masses are smaller than
$E_{\mathrm{cm}}/2$.  When $E_{\mathrm{cm}}/2$ is in the vicinity of one of
the quark masses, the strong interactions cause large deviations from this
formula.  The most dramatic such effect is the appearance of bound states just
below $E_{\mathrm{cm}} = 2m_q$, manifested as very sharp spikes in the cross
section.

Experimental measurements of the cross section for $e^{+}e^{-}$ annihilation
to hadrons between 2.5 and 40 GeV are shown in Fig.~\ref{fig:5.3}.  The data
shows three distinct regions: a low-energy region in which $u$, $d$, and $s$
quark pairs are produced; a region above the threshold for production of $c$
quark pairs; and a region also above the threshold for $b$ quark pairs.  The
prediction \eqref{eq:5.16} is shown as a set of solid lines; it agrees quite
well with the data in each region, as long as the energy is well away from the
thresholds where the high-energy approximation breaks down.  The dotted
curves show an improved theoretical prediction, including higher-order
corrections from QCD, which we will discuss in Section~17.2.  This explanation
of the $e^{+}e^{-}$ annihilation cross section is a remarkable success of QCD.
In particular, experimental verification of the factor of 3 in
\eqref{eq:5.16} is one piece of evidence for the existence of color.

\begin{figure}[t]
  \centering
  \includegraphics[width=0.7\textwidth]{fig5.3.png}
  \caption{Experimental measurements of the total cross section for the
    reaction $e^{+}e^{-}\to$ hadrons, from the data compilation of M.~Swartz,
    \emph{Phys.\ Rev.\ D} \textbf{53}, 5268 (1996).  Complete references to the
    various experiments are given there.  The measurements are compared to
    theoretical predictions from Quantum Chromodynamics, as explained in the
    text.  The solid line is the simple prediction \eqref{eq:5.16}.}
  \label{fig:5.3}
\end{figure}

The angular dependence of the differential cross section is also observed
experimentally.%
\footnote{The basic features of hadron production in high-energy $e^{+}e^{-}$
annihilation are reviewed by P.~Duinker, \emph{Rev.\ Mod.\ Phys.} \textbf{54},
325 (1982).}  At high energy the hadrons appear in jets, clusters of several
hadrons all moving in approximately the same direction.  In most cases there
are two jets, with back-to-back momenta, and these indeed have the angular
dependence $(1 + \cos^2\theta)$.

\section{$e^{+}e^{-}\to\mu^{+}\mu^{-}$: Helicity Structure}

The unpolarized cross section for a reaction is generally easy to calculate
(and to measure) but hard to understand.  Where does the
$(1 + \cos^2\theta)$ angular dependence come from?  We can answer this
question by computing the $e^{+}e^{-}\to\mu^{+}\mu^{-}$ cross section for each
set of spin orientations separately.

First we must choose a basis of polarization states.  To get a simple answer
in the high-energy limit, the best choice is to quantize each spin along the
direction of the particle's motion, that is, to use states of definite
helicity.  Recall that in the massless limit, the left- and right-handed
helicity states of a Dirac particle live in different representations of the
Lorentz group.  We might therefore expect them to behave independently, and in
fact they do.

In this section we will compute the polarized $e^{+}e^{-}\to\mu^{+}\mu^{-}$
cross sections, using the helicity basis, in two different ways: first, by
using trace technology but with the addition of helicity projection operators
to project out the desired left- or right-handed spinors; and second, by
plugging explicit expressions for these spinors directly into our formula for
the amplitude $\mathcal{M}$.  Throughout this section we work in the
high-energy limit where all fermions are effectively massless.  (The
calculation can be done for lower energy, but it is much more difficult and no
more instructive.)%
\footnote{A general treatment of this calculation, retaining all masses, is
given by F.~M.~Renard, \emph{Basics of Electron-Positron Collisions} (Editions
Fronti\`eres, 1981).}

Our starting point for both methods of calculating the polarized cross section
is the amplitude
\begin{equation}
  \mathcal{M} = \frac{e^2}{q^2}\
    \bigl( \bar{v}(p')\gamma^\mu u(p) \bigr)\
    \bigl( \bar{u}(k)\gamma_\mu v(k') \bigr) .
  \label{eq:5.19}
\end{equation}
We use the helicity projection operators
\begin{equation}
  P_R = \frac{1+\gamma^5}{2},\qquad P_L = \frac{1-\gamma^5}{2} .
\end{equation}
The left- and right-handed spinors are defined by
\begin{equation}
  u_L = P_L u, \qquad u_R = P_R u, \qquad v_L = P_L v, \qquad v_R = P_R v .
\end{equation}
To compute the amplitude for definite helicities, we insert the projection
operators before each spinor and use the trace technology.  The result, after
summing over the (now definite) helicity states, gives the polarized
cross-section formula.  For example, for the process
$e^{-}_R e^{+}_L \to \mu^{-}_R \mu^{+}_L$, one finds
\begin{equation}
  \Bigl( \frac{d\sigma}{d\Omega} \Bigr)(e^{-}_R e^{+}_L \to \mu^{-}_R
    \mu^{+}_L)
    = \frac{\alpha^2}{4E_{\mathrm{cm}}^2}\ (1 + \cos\theta)^2 .
  \label{eq:5.21}
\end{equation}
Note that the cross section \eqref{eq:5.21} vanishes at $\theta = 180^{\circ}$.
This is just what we would expect, since for $\theta = 180^{\circ}$, the total
angular momentum of the final state is opposite to that of the initial state
(see Fig.~\ref{fig:5.4}).

\begin{figure}[t]
  \centering
  \includegraphics[width=0.5\textwidth]{fig5.4.png}
  \caption{Conservation of angular momentum requires that if the $z$-component
    of angular momentum is measured, it must have the same value as
    initially.}
  \label{fig:5.4}
\end{figure}

This completes our first calculation of the polarized
$e^{+}e^{-}\to\mu^{+}\mu^{-}$ cross sections.  We will now redo the
calculation in a manner that is more straightforward, more enlightening, and
no more difficult.  We will calculate the amplitude $\mathcal{M}$ (rather than
the squared amplitude) directly, using explicit values for the spinors and
$\gamma$ matrices.  This method does have its drawbacks: It forces us to
specialize to a particular frame of reference much sooner, so manifest Lorentz
invariance is lost.  More pragmatically, it is very cumbersome except in the
nonrelativistic and ultra-relativistic limits.

Consider again the amplitude
\begin{equation}
  \mathcal{M} = \frac{e^2}{q^2}\
    \bigl( \bar{v}(p')\gamma^\mu u(p) \bigr)\
    \bigl( \bar{u}(k)\gamma_\mu v(k') \bigr) .
  \label{eq:5.25}
\end{equation}
In the high-energy limit, our general expressions for Dirac spinors become
\begin{equation}
  u(p) = \begin{pmatrix}
    \sqrt{p\cdot\sigma}\,\xi\\
    \sqrt{p\cdot\bar{\sigma}}\,\xi
  \end{pmatrix}
  \xrightarrow{E\to\infty}
  \sqrt{E}\begin{pmatrix}
    \frac{1+\hat{\mathbf{p}}\cdot\bm{\sigma}}{2}\xi\\
    \frac{1+\hat{\mathbf{p}}\cdot\bm{\sigma}}{2}\xi
  \end{pmatrix};\qquad
  v(p) = \begin{pmatrix}
    \sqrt{p\cdot\sigma}\,\xi\\
    -\sqrt{p\cdot\bar{\sigma}}\,\xi
  \end{pmatrix}
  \xrightarrow{E\to\infty}
  \sqrt{E}\begin{pmatrix}
    \frac{1-\hat{\mathbf{p}}\cdot\bm{\sigma}}{2}\xi\\
    -\frac{1-\hat{\mathbf{p}}\cdot\bm{\sigma}}{2}\xi
  \end{pmatrix} .
  \label{eq:5.26}
\end{equation}
A right-handed spinor satisfies $(\hat{\mathbf{p}}\cdot\bm{\sigma})\xi = +\xi$,
while a left-handed spinor has $(\hat{\mathbf{p}}\cdot\bm{\sigma})\xi = -\xi$.
(Remember once again that for antiparticles, the handedness of the spinor is
the opposite of the handedness of the particle.)  We must evaluate expressions
of the form $\bar{v}\gamma^\mu u$, so we need the identity
\begin{equation}
  P_R\gamma^\mu P_L = \gamma^\mu P_L,\qquad
  P_L\gamma^\mu P_R = \gamma^\mu P_R .
\end{equation}
Thus we see explicitly that the amplitude is zero when one of the spinors is
left-handed and the other is right-handed.  In the language of Chapter~1, the
Clebsch--Gordan coefficients that couple the vector photon to the product of
such spinors are zero; those coefficients are just the off-block-diagonal
elements of the matrix $\gamma^0\gamma^\mu$ (in the chiral representation).

Let us choose $p$ and $p'$ to be in the $\hat{z}$ directions, and first
consider the case where the electron is right-handed and the positron is
left-handed.  Thus for the electron we have $\xi = \binom{1}{0}$,
corresponding to spin up in the $\hat{z}$ direction, while for the positron we
have $\xi = \binom{1}{0}$, also corresponding to (physical) spin up in the
$\hat{z}$-direction.  Both particles have $(\hat{\mathbf{p}}\cdot\bm{\sigma})
\xi = +\xi$, so the spinors are
\begin{equation}
  u(p) = \sqrt{2E}\begin{pmatrix}\binom{1}{0}\\ \binom{1}{0}\end{pmatrix},
  \qquad v(p') = \sqrt{2E}\begin{pmatrix}\binom{0}{1}\\ -\binom{0}{1}\end{pmatrix}.
  \label{eq:5.28}
\end{equation}
The electron half of the matrix element is therefore
\begin{equation}
  \bar{v}(p')\gamma^\mu u(p) = 2E(0, -i)\cdots = -2E(0, 1, i, 0) .
  \label{eq:5.29}
\end{equation}
We can interpret this expression by saying that the virtual photon has
circular polarization in the $+z$-direction; its polarization vector is
$\epsilon^{+} = (1/\sqrt{2})(\hat{x}+i\hat{y})$.

Next we must calculate the muon half of the matrix element.  Let the
$\mu^{-}$ be emitted at an angle $\theta$ to the $\hat{z}$-axis, and consider
first the case where it is right-handed (and the $\mu^{+}$ is therefore
left-handed).  To calculate $\bar{u}(k)\gamma^\mu v(k')$ we could go back to
expressions \eqref{eq:5.26}, but then it would be necessary to find the
correct spinors $\xi$ corresponding to polarization along the muon momentum.
It is much easier to use a trick: Since any expression of the form
$\bar{u}\gamma^\mu v$ transforms like a 4-vector, we can just rotate the result
\eqref{eq:5.29}.  Rotating that vector by an angle $\theta$ in the $xz$-plane,
we find
\begin{equation}
  \bar{u}(k)\gamma^\mu v(k') = \bigl[ \bar{v}(k')\gamma^\mu u(k) \bigr]^{*}
    = \bigl[ -2E(0, \cos\theta, i, -\sin\theta) \bigr]^{*}
    = -2E(0, \cos\theta, -i, -\sin\theta) .
  \label{eq:5.30}
\end{equation}
This vector can also be interpreted as the polarization of the virtual photon;
when it has a nonzero overlap with \eqref{eq:5.29}, we get a nonzero
amplitude.  Plugging \eqref{eq:5.29} and \eqref{eq:5.30} into
\eqref{eq:5.25}, we see that the amplitude is
\begin{equation}
  \mathcal{M}(e^{-}_R e^{+}_L \to \mu^{-}_R\mu^{+}_L)
    = \frac{e^2}{q^2}(2E)^2(-\cos\theta - 1) = -e^2(1+\cos\theta),
  \label{eq:5.31}
\end{equation}
in agreement with \eqref{eq:1.6}, and also with \eqref{eq:5.21}.  The
differential cross section for this set of helicities can now be obtained in
the same way as above, yielding \eqref{eq:5.21}.

We can calculate the other three nonvanishing helicity amplitudes in an
analogous manner.  For a left-handed electron and a right-handed positron, we
easily find
\begin{equation}
  \bar{v}(p')\gamma^\mu u(p) = -2E(0, 1, -i, 0) = -2E\sqrt{2}\,\epsilon^{\mu}_{-} .
\end{equation}
Perform a rotation to get the vector corresponding to a left-handed
$\mu^{-}$ and a right-handed $\mu^{+}$:
\begin{equation}
  \bar{u}(k)\gamma^\mu v(k') = -2E(0, \cos\theta, i, -\sin\theta) .
\end{equation}
Putting the pieces together in various ways yields the remaining amplitudes,
\begin{equation}
  \begin{aligned}
    \mathcal{M}(e^{-}_L e^{+}_R \to \mu^{-}_L\mu^{+}_R) &= -e^2(1+\cos\theta);\\
    \mathcal{M}(e^{-}_L e^{+}_R \to \mu^{-}_R\mu^{+}_L) &=
    \mathcal{M}(e^{-}_R e^{+}_L \to \mu^{-}_L\mu^{+}_R) = -e^2(1-\cos\theta).
  \end{aligned}
  \label{eq:5.32}
\end{equation}

\section{$e^{+}e^{-}\to\mu^{+}\mu^{-}$: Nonrelativistic Limit}

Now let us go to the other end of the energy spectrum, and discuss the
reaction $e^{+}e^{-}\to\mu^{+}\mu^{-}$ in the extreme nonrelativistic limit.
When $E$ is barely larger than $m_\mu$, our previous result \eqref{eq:5.12}
for the unpolarized differential cross section becomes
\begin{equation}
  \frac{d\sigma}{d\Omega} \xrightarrow{|\mathbf{k}|\to 0}
    \frac{\alpha^2}{4E^2}\ \frac{|\mathbf{k}|}{E}\ \frac{1}{1}
    = \frac{\alpha^2|\mathbf{k}|}{4E^3} .
  \label{eq:5.33}
\end{equation}
We can recover this result, and also learn something about the spin dependence
of the reaction, by evaluating the amplitude with explicit spinors.

Once again we begin with the matrix element
\begin{equation}
  \mathcal{M} = \frac{e^2}{q^2}\ \bar{v}(p')\gamma^\mu u(p)\
    \bar{u}(k)\gamma_\mu v(k') .
\end{equation}
The electron and positron are still very relativistic, so this expression will
be simplest if we choose them to have definite helicity.  Let the electron be
right-handed, moving in the $+\hat{z}$-direction, and the positron be
left-handed, moving in the $-\hat{z}$-direction.  Then from Eq.~\eqref{eq:5.29}
we have
\begin{equation}
  \bar{v}(p')\gamma^\mu u(p) = -2E(0, 1, i, 0) .
  \label{eq:5.34}
\end{equation}
In the other half of the matrix element we should use the nonrelativistic
expressions
\begin{equation}
  u(k) = \sqrt{m}\begin{pmatrix}\xi\\ \xi\end{pmatrix},\qquad
  v(k') = \sqrt{m}\begin{pmatrix}\xi'\\ -\xi'\end{pmatrix} .
  \label{eq:5.35}
\end{equation}
Keep in mind, in the discussion of this section, that the spinor $\xi'$ gives
the flipped spin of the antiparticle.  Leaving the muon spinors $\xi$ and
$\xi'$ undetermined for now, we can easily compute
\begin{equation}
  \bar{u}(k)\gamma^\mu v(k') = \cdots = -m\bigl( \xi^{\dagger}\xi',
    -\xi^{\dagger}\bm{\sigma}\xi' \bigr) .
  \label{eq:5.36}
\end{equation}
To evaluate $\mathcal{M}$, we simply dot \eqref{eq:5.34} into \eqref{eq:5.36}
and multiply by $e^2/q^2 = e^2/4m^2$.  The result is
\begin{equation}
  \mathcal{M}(e^{-}_R e^{+}_L \to \mu^{+}\mu^{-})
    = -2e^2\,\xi^{\dagger}\begin{pmatrix}0&0\\0&1\end{pmatrix}\xi' .
  \label{eq:5.37}
\end{equation}
Since there is no angular dependence in this expression, the muons are equally
likely to come out in any direction.  More precisely, they are emitted in an
$S$-wave; their orbital angular momentum is zero.  Angular momentum
conservation therefore requires that the total spin of the final state equal
1, and indeed the matrix product gives zero unless both the muon and the
antimuon have spin up along the $\hat{z}$-axis (see Fig.~\ref{fig:5.5}).

\begin{figure}[t]
  \centering
  \includegraphics[width=0.5\textwidth]{fig5.5.png}
  \caption{In the nonrelativistic limit the total spin of the system is
    conserved, and thus the muons are produced with both spins up along the
    $\hat{z}$-axis.}
  \label{fig:5.5}
\end{figure}

To find the total rate for this process, we sum over muon spins to obtain
$|\mathcal{M}|^2 = 4e^4$, which yields the cross section
\begin{equation}
  \frac{d\sigma}{d\Omega}(e^{-}_R e^{+}_L \to \mu^{+}\mu^{-})
    = \frac{\alpha^2|\mathbf{k}|}{4m_\mu^2 E^2} .
  \label{eq:5.38}
\end{equation}
The same expression holds for a left-handed electron and a right-handed
positron.  Thus the spin-averaged cross section is just
$2\cdot(1/4)$ times this expression, in agreement with \eqref{eq:5.33}.

\subsection{Bound States}

Until now we have considered the initial and final states of scattering
processes to be states of isolated single particles.  Very close to threshold,
however, the Coulomb attraction of the muons should become an important
effect.  Just below threshold, we can still form $\mu^{+}\mu^{-}$ pairs in
electromagnetic bound states.

The treatment of bound states in quantum field theory is a rich and complex
subject, but one that lies mainly beyond the scope of this book.%
\footnote{See, for example, the review by C.~Itzykson and J.-B.~Zuber, \emph{Quantum
Field Theory} (McGraw-Hill, 1980), and the references therein.}  Fortunately,
many of the familiar bound systems in Nature can be treated (at least to a
good first approximation) as nonrelativistic systems, in which the internal
motions are slow.  The process of creating the constituent particles out of
the vacuum is still a relativistic effect, requiring quantum field theory for
its proper description.  In this section we will develop a formalism for
computing the amplitudes for creation and annihilation of two-particle,
nonrelativistic bound states.  We begin with a computation of the cross
section for producing a $\mu^{+}\mu^{-}$ bound state in $e^{+}e^{-}$
annihilation.

Consider first the case where the spins of the electron and positron both
point in the same direction.  The general bound state of a fermion and its
antifermion has total spin $S = 0$ or $S = 1$, and a spatial wavefunction
$\psi(\mathbf{r})$.  The nonrelativistic limit of the amplitude for producing
a bound state $B$ is proportional to the wavefunction at the origin,
\begin{equation}
  \mathcal{M}(e^{+}e^{-}\to B) \propto \psi(0) .
  \label{eq:5.50}
\end{equation}
Carrying out the computation (see the details in the text), we obtain
\begin{equation}
  \mathcal{M}(e^{+}e^{-}\to B) = 2e^2\sqrt{2M}\,\psi(0)\,
    \epsilon(\mathbf{n})\cdot\hat{z}
  \label{eq:5.51}
\end{equation}
for the case of a spin-1 bound state with polarization vector
$\bm{\epsilon}$, where $M = 2m_\mu$ is the bound-state mass and
$\mathbf{n}$ is the direction of the muon momentum.  Assuming an unpolarized
electron beam, and summing Eq.~(5.52) over the three possible directions
of $\mathbf{n}$, we find the following expression for the total cross section
for production of the bound state:
\begin{equation}
  \sigma(e^{+}e^{-}\to B) =
    \frac{(2\pi)^4}{2E^2\,2E}\ \frac{1}{2}\ \delta(P^0 - M^0)\
    \frac{1}{3}(4e^4)|\psi(0)|^2 .
  \label{eq:5.53}
\end{equation}
Notice that the 1-body phase space integral can remove only three of the four
delta functions.  It is conventional to rewrite the last delta function using
$\delta(P^0 - K^0) = 2K^0\,\delta(P^2 - K^2)$.  Then
\begin{equation}
  \sigma(e^{+}e^{-}\to B) = 64\pi^3\alpha^2\ \frac{|\psi(0)|^2}{M}\
    \delta(E^2_{\mathrm{cm}} - M^2) .
  \label{eq:5.54}
\end{equation}
The last delta function enforces the constraint that the total center-of-mass
energy must equal the bound-state mass; thus, the bound state is produced as a
resonance in $e^{+}e^{-}$ annihilation.  If the bound state has a finite
lifetime, this delta function will be broadened into a resonance peak.  In
practice, the intrinsic spread of the $e^{+}e^{-}$ beam energy is often a more
important broadening mechanism.  In either case, \eqref{eq:5.54} correctly
predicts the area under the resonance peak.

If the bound state $B$ can be produced from $e^{+}e^{-}$, it can also annihilate
back to $e^{+}e^{-}$, or to any other sufficiently light lepton pair.
According to Eq.~(4.86), the total width for this decay mode is given by
\begin{equation}
  \Gamma(B \to e^{+}e^{-}) = \frac{1}{2M}\ \int d\Pi_2\ |\mathcal{M}|^2 ,
  \label{eq:5.55}
\end{equation}
where $\mathcal{M}$ is just the complex conjugate of the matrix element
\eqref{eq:5.51} we used to compute $B$ production.  Thus
\begin{equation}
  \Gamma(B \to e^{+}e^{-}) = 16\pi\alpha^2\ \frac{|\psi(0)|^2}{M^2}\
    \cdot\ 3\ \cdot\ \cdots .
  \label{eq:5.57}
\end{equation}
This formula is summed over the possible final electron polarization states.
It is easiest to evaluate by averaging over the three possible values of
$\mathbf{n}$.  We thus obtain \eqref{eq:5.57}.  The formula for the decay
width of $B$ is very similar to that for the production cross section, and
this is no surprise: Both calculations involve the square of the same matrix
element, summed over initial and final polarizations.  The two calculations
differed only in how we formed the polarization averages, and in the
phase-space factors.  By this logic, the relation we have found between the
two quantities,
\begin{equation}
  \sigma(e^{+}e^{-}\to B) = 4\pi^2\cdot\frac{3}{M^2}\cdot
    \Gamma(B\to e^{+}e^{-})\ \delta(E^2_{\mathrm{cm}} - M^2),
  \label{eq:5.58}
\end{equation}
is very general and completely independent of the details of the matrix
element computation.  The factor 3 in \eqref{eq:5.58} came from the
orientation average for $\mathbf{n}$; for a spin-$J$ bound state, this factor
would be $(2J+1)$.

The most famous application of this formalism is to bound states not of muons
but of quarks: quarkonium.  We saw the experimental evidence for $q\bar{q}$
bound states (the $J/\psi$ and $\Upsilon$, for example) in Fig.~\ref{fig:5.3}.
(The resonance peaks are much too high and too narrow to show in the figure,
but their sizes have been carefully measured.)  Equations \eqref{eq:5.54} and
\eqref{eq:5.57} must be multiplied by a color factor of 3 to give the
production cross section and decay width for a spin-1 $q\bar{q}$ bound state.
The value $\psi(0)$ of the $q\bar{q}$ wavefunction at the origin cannot be
computed from first principles, but can be estimated from a nonrelativistic
model of the $q\bar{q}$ spectrum with a phenomenologically chosen potential.
Alternatively, we can use the formula
\begin{equation}
  \Gamma\bigl( B(q\bar{q}) \to e^{+}e^{-} \bigr) = 16\pi\alpha^2 Q^2
    \frac{|\psi(0)|^2}{M^2}
  \label{eq:5.59}
\end{equation}
to measure $\psi(0)$ for a $q\bar{q}$ bound state.  For example, the $1S$
spin-1 state of $s\bar{s}$, the $\phi$ meson, has an $e^{+}e^{-}$ partial width
of 1.4 keV and a mass of 1.02 GeV.  From this we can infer
$|\psi(0)|^2 = (1.2\ \mathrm{fm})^{-3}$.  This result is physically
reasonable, since hadronic dimensions are typically $\sim 1$ fm.

Our viewpoint in this section has been quite different from that of earlier
sections: Instead of computing everything from first principles, we have
pieced together an approximate formula using a bit of quantum field theory and
a bit of nonrelativistic quantum mechanics.  In principle, however, we could
treat bound states entirely in the relativistic formalism.  Consider the
annihilation of an $e^{+}e^{-}$ pair to form a $\mu^{+}\mu^{-}$ bound state,
which subsequently decays back into $e^{+}e^{-}$.  In our present formalism we
might represent this process by the diagram
\begin{equation}
  e^{+}e^{-} \to B \to e^{+}e^{-} .
\end{equation}
The net process is simply $e^{+}e^{-}\to e^{+}e^{-}$ (Bhabha scattering).
What would happen if we tried to compute the Bhabha scattering cross section
directly in QED perturbation theory?  Obviously there is no
$\mu^{+}\mu^{-}$ contribution in the tree-level diagrams.  As we go to higher
orders in the perturbation series, however, we find (among others) the
following set of diagrams, in which a $\mu^{+}\mu^{-}$ loop is exchanged
between the two electron lines.  At most values of $E_{\mathrm{cm}}$, these
diagrams give only a small correction to the tree-level expression.  But when
$E_{\mathrm{cm}}$ is near the $\mu^{+}\mu^{-}$ threshold, the diagrams involving
the exchange of photons within the muon loop contain the Coulomb interaction
between the muons, and therefore become quite large.  One must sum over all
such diagrams, and it can be shown that this summation is equivalent to
solving the nonrelativistic Schr\"odinger equation.%
\footnote{This analysis is carried out in Berestetskii, Lifshitz, and
Pitaevskii (1982).}  The final prediction is that the cross section contains a
resonance peak, whose area is given by \eqref{eq:5.54} and whose width is
given by \eqref{eq:5.57}.

\section{Crossing Symmetry}

\subsection{Electron-Muon Scattering}

Now that we have completed our discussion of the process
$e^{+}e^{-}\to\mu^{+}\mu^{-}$, let us consider a different but closely related
QED process: electron-muon scattering, or $e^{-}\mu^{-}\to e^{-}\mu^{-}$.  The
lowest-order Feynman diagram is just the previous one turned on its side:
\begin{equation}
  i\mathcal{M} = \frac{(-ie)^2}{q^2}\ \bar{u}(p_1')\gamma^\mu u(p_1)\
    \bar{u}(p_2')\gamma_\mu u(p_2) .
\end{equation}
The relation between the processes $e^{+}e^{-}\to\mu^{+}\mu^{-}$ and
$e^{-}\mu^{-}\to e^{-}\mu^{-}$ becomes clear when we compute the squared
amplitude, averaged and summed over spins:
\begin{equation}
  \frac{1}{4}\sum_{\text{spins}}|\mathcal{M}|^2
    = \frac{e^4}{4q^4}\ \mathrm{tr}\bigl[
      (\slashed{p}_1' + m_e)\gamma^\mu(\slashed{p}_1 + m_e)\gamma^\nu \bigr]\
      \mathrm{tr}\bigl[ (\slashed{p}_2' + m_\mu)\gamma_\mu(\slashed{p}_2
        + m_\mu)\gamma_\nu \bigr] .
\end{equation}
This is exactly the same as our result \eqref{eq:5.4} for
$e^{+}e^{-}\to\mu^{+}\mu^{-}$, with the replacements
\begin{equation}
  p \to p_1,\qquad p' \to -p_1',\qquad k \to p_2',\qquad k' \to -p_2 .
  \label{eq:5.60}
\end{equation}
So instead of evaluating the traces from scratch, we can just make the same
replacements in our previous result, Eq.~\eqref{eq:5.10}.  Setting $m_e = 0$,
we find
\begin{equation}
  \frac{1}{4}\sum_{\text{spins}}|\mathcal{M}|^2
    = \frac{e^4}{4q^4}\ \Bigl[ (p_1\cdot p_2)(p_1'\cdot p_2')
      + (p_1\cdot p_2')(p_1'\cdot p_2) - m_\mu^2(p_1\cdot p_1') \Bigr] .
  \label{eq:5.61}
\end{equation}
To evaluate this expression, we must work out the kinematics, which will be
completely different.  Working in the center-of-mass frame, we make the
following assignments:
\begin{equation}
  \begin{aligned}
    p_1 &= (k, \mathbf{k}),\qquad p_1' = (k, \mathbf{k}');\qquad
    p_2 = (E, -\mathbf{k}),\qquad p_2' = (E, -\mathbf{k}').
  \end{aligned}
\end{equation}
The combinations we need are
\begin{equation}
  \begin{aligned}
    p_1\cdot p_2 = p_1'\cdot p_2' &= k(E+k);\\
    p_1\cdot p_1' &= k^2(1-\cos\theta);\\
    p_1\cdot p_2' = p_1'\cdot p_2 &= k(E+k\cos\theta);\\
    q^2 &= -2p_1\cdot p_1' = -2k^2(1-\cos\theta) .
  \end{aligned}
\end{equation}
Our expression for the squared matrix element now becomes
\begin{equation}
  \frac{1}{4}\sum_{\text{spins}}|\mathcal{M}|^2
    = \frac{e^4}{4k^2E^2(1-\cos\theta)^2}\ \Bigl[ k^2(E+k)^2
      + k^2(E+k\cos\theta)^2 - m_\mu^2 k^2(1-\cos\theta) \Bigr] .
  \label{eq:5.62}
\end{equation}
To find the cross section from this expression, we use Eq.~(4.84),
which in the case where one particle is massless takes the simple form
\begin{equation}
  \Bigl( \frac{d\sigma}{d\Omega} \Bigr)_{\mathrm{CM}}
    = \frac{|\mathcal{M}|^2}{64\pi^2(E+k)^2} .
  \label{eq:5.63}
\end{equation}
Thus we have our result for unpolarized electron-muon scattering in the
center-of-mass frame:
\begin{equation}
  \Bigl( \frac{d\sigma}{d\Omega} \Bigr)_{\mathrm{CM}}
    = \frac{e^4}{64\pi^2(E+k)^2}\ \frac{1}{4k^2E^2(1-\cos\theta)^2}\
      \Bigl[ \cdots \Bigr] ,
  \label{eq:5.64}
\end{equation}
where $k = \sqrt{E^2-m_\mu^2}$.  In the high-energy limit where we can set
$m_\mu = 0$, the differential cross section becomes
\begin{equation}
  \frac{d\sigma}{d\Omega} \xrightarrow{E\gg m_\mu}
    \frac{\alpha^2}{4E_{\mathrm{cm}}^2}\ \frac{1}{\sin^4(\theta/2)}
      \Bigl( \frac{1+\cos^2\theta}{4} \Bigr) .
  \label{eq:5.65}
\end{equation}

\subsection{Crossing Symmetry and the Mandelstam Variables}

The manipulations of the previous subsection are a special case of a more
general relation between scattering amplitudes, called \emph{crossing
symmetry}.  Consider a general two-body scattering process.  It is
conventional to introduce the Mandelstam variables
\begin{equation}
  s = (p+p')^2,\qquad t = (k-p)^2,\qquad u = (k'-p)^2 ,
  \label{eq:5.69}
\end{equation}
where $p$ and $p'$ are the initial momenta and $k$ and $k'$ are the final
momenta.  The three variables satisfy
\begin{equation}
  s + t + u = \sum_{i=1}^{4} m_i^2 ,
  \label{eq:5.73}
\end{equation}
where the sum runs over the four external particles.  This identity is easy to
prove by adding up the terms on the right-hand side of Eqs.~\eqref{eq:5.69},
and applying momentum conservation in the form
$(p+p'-k-k')^2 = 0$.

The crossing symmetry relates the amplitude for a process to the amplitude
for the related process in which some initial particles are moved to the final
state (replacing their momenta by their negatives).  For example, the
amplitude for $e^{+}e^{-}\to\mu^{+}\mu^{-}$ in terms of the Mandelstam
variables can be obtained from the amplitude for $e^{-}\mu^{-}\to e^{-}\mu^{-}$
by the substitutions
\begin{equation}
  p \leftrightarrow k,\qquad p' \to -p',\qquad k' \to -k' .
  \label{eq:5.68}
\end{equation}
In the center-of-mass frame with all four masses equal to $m$, we have
\begin{equation}
  s = (p+p')^2 = (2E)^2 = E^2_{\mathrm{cm}},\qquad
  t = (k-p)^2 = -2p^2(1-\cos\theta),\qquad
  u = (k'-p)^2 = -2p^2(1+\cos\theta) .
  \label{eq:5.72}
\end{equation}
Thus we see that $t\to 0$ as $\theta\to 0$, while $u\to 0$ as $\theta\to\pi$.
(When the masses are not all equal, the limiting values of $t$ and $u$ will
shift slightly.)  Note from \eqref{eq:5.72} that when all four particles have
mass $m$, the sum of the Mandelstam variables is
$s+t+u = 4E^2 - 4p^2 = 4m^2$.  This is a special case of the more general
relation \eqref{eq:5.73}.

\section{Compton Scattering}

We now move on to consider a somewhat different QED process: Compton
scattering, or $e^{-}\gamma\to e^{-}\gamma$.  We will calculate the unpolarized
cross section for this reaction, to lowest order in $\alpha$.  The calculation
will employ all the machinery we have developed so far, including the
Mandelstam variables of the previous section.  We will also develop some new
technology for dealing with external photons.

This is our first example of a calculation involving two diagrams:
\begin{equation}
  i\mathcal{M} = \bar{u}(p')(-ie\gamma^\mu)\epsilon^{*}_{\mu}(k')\
    \frac{i(\slashed{p}+\slashed{k}+m)}{(p+k)^2-m^2}\
    (-ie\gamma^\nu)\epsilon_\nu(k)u(p)
    + \bar{u}(p')(-ie\gamma^\mu)\epsilon_\nu(k)\
      \frac{i(\slashed{p}-\slashed{k}'+m)}{(p-k')^2-m^2}\
      (-ie\gamma^\nu)\epsilon^{*}_\mu(k')u(p) .
\end{equation}
As usual, the Feynman rules tell us exactly how to write down an expression
for $\mathcal{M}$.  Note that since the fermion portions of the two diagrams
are identical, there is no relative minus sign between the two terms.  Using
$\epsilon_\mu(k)$ and $\epsilon^{*}_\mu(k')$ to denote the polarization vectors
of the initial and final photons, we have
\begin{equation}
  \begin{aligned}
    i\mathcal{M} &= -ie^2\ \epsilon^{*}_{\nu}(k')\epsilon_\mu(k)\
      \bar{u}(p')\ \Bigl[
        \frac{\gamma^\mu(\slashed{p}+\slashed{k}+m)\gamma^\nu}{(p+k)^2-m^2}
        + \frac{\gamma^\nu(\slashed{p}-\slashed{k}'+m)\gamma^\mu}
          {(p-k')^2-m^2} \Bigr]\ u(p) .
  \end{aligned}
  \label{eq:5.74}
\end{equation}
We can make a few simplifications before squaring this expression.  Since
$p^2 = m^2$ and $k^2 = 0$, the denominators of the propagators are
\begin{equation}
  (p+k)^2 - m^2 = 2p\cdot k
  \qquad\text{and}\qquad
  (p-k')^2 - m^2 = -2p'\cdot k' .
\end{equation}
To simplify the numerators, we use a bit of Dirac algebra:
\begin{equation}
  (\slashed{p}+m)\gamma^\nu u(p) = (2p^\nu - \gamma^\nu\slashed{p}
    + \gamma^\nu m)u(p) = 2p^\nu u(p) .
\end{equation}
Using this trick on the numerator of each propagator, we obtain
\begin{equation}
  i\mathcal{M} = -ie^2\ \epsilon^{*}_{\nu}(k')\epsilon_\mu(k)\
    \bar{u}(p')\ \Bigl[ \frac{\gamma^\mu\slashed{k}\gamma^\nu}{2p\cdot k}
      - \frac{\gamma^\nu\slashed{k}'\gamma^\mu}{2p'\cdot k'} \Bigr]\ u(p) .
  \label{eq:5.74b}
\end{equation}

\subsection{Photon Polarization Sums}

The next step in the calculation will be to square this expression for
$\mathcal{M}$ and sum (or average) over electron and photon polarization
states.  The sum over electron polarizations can be performed as before, using
the identity $\sum u(p)\bar{u}(p) = \slashed{p} + m$.  Fortunately, there is a
similar trick for summing over photon polarization vectors.  The correct
prescription is to make the replacement
\begin{equation}
  \sum_{\text{polarizations}}\ \epsilon^{*}_\mu(k)\epsilon_\nu(k)
    \to -g_{\mu\nu} .
  \label{eq:5.75}
\end{equation}
The arrow indicates that this is not an actual equality.  Nevertheless, the
replacement is valid as long as both sides are dotted into the rest of the
expression for a QED amplitude $\mathcal{M}$.

To derive this formula, let us consider an arbitrary QED process involving an
external photon with momentum $k$:
\begin{equation}
  \mathcal{M}(k) = i\mathcal{M}^{\mu}(k)\epsilon^{*}_\mu(k) .
  \label{eq:5.76}
\end{equation}
Since the amplitude always contains $\epsilon^{*}_\mu(k)$, we have extracted
this factor and defined $\mathcal{M}^{\mu}(k)$ to be the rest of the amplitude
$\mathcal{M}$.  The cross section will be proportional to
\begin{equation}
  \sum_{\epsilon} \bigl| \epsilon^{*}_\mu(k)\mathcal{M}^{\mu}(k) \bigr|^2
    = \sum_{\epsilon} \epsilon^{*}_\mu(k)\epsilon_\nu(k)\,
      \mathcal{M}^{\mu}(k)\mathcal{M}^{\nu*}(k) .
\end{equation}
For simplicity, we orient $k$ in the 3-direction: $k^\mu = (k, 0, 0, k)$.
Then the two transverse polarization vectors, over which we are summing, can
be chosen to be
\begin{equation}
  \epsilon_1 = (0, 1, 0, 0);\qquad \epsilon_2 = (0, 0, 1, 0) .
\end{equation}
With these conventions, we have
\begin{equation}
  \sum_{\epsilon} \bigl| \epsilon^{*}_\mu(k)\mathcal{M}^{\mu}(k) \bigr|^2
    = |\mathcal{M}^1(k)|^2 + |\mathcal{M}^2(k)|^2 .
  \label{eq:5.77}
\end{equation}
Now recall from Chapter~4 that external photons are created by the interaction
term $\int d^4x\ e j^\mu A_\mu$, where $j^\mu = \bar{\psi}\gamma^\mu\psi$ is
the Dirac vector current.  Therefore we expect $\mathcal{M}^{\mu}(k)$ to be
given by a matrix element of the Heisenberg field $j^\mu$:
\begin{equation}
  \mathcal{M}^{\mu}(k) = \int d^4x\ e^{ik\cdot x}\ \langle f|j^\mu(x)|i\rangle ,
  \label{eq:5.78}
\end{equation}
where the initial and final states include all particles except the photon in
question.

From the classical equations of motion, we know that the current $j^\mu$ is
conserved: $\partial_\mu j^\mu(x) = 0$.  Provided that this property still
holds in the quantum theory, we can dot $k_\mu$ into \eqref{eq:5.78} to
obtain
\begin{equation}
  k_\mu\mathcal{M}^{\mu}(k) = 0 .
  \label{eq:5.79}
\end{equation}
The amplitude $\mathcal{M}$ vanishes when the polarization vector
$\epsilon^\mu(k)$ is replaced by $k^\mu$.  This famous relation is known as
the \emph{Ward identity}.  It is essentially a statement of current
conservation, which is a consequence of the gauge symmetry \eqref{eq:4.6} of
QED.  We will give a formal proof of the Ward identity in Section~7.4, and
discuss a number of subtle points skimmed over in this quick ``derivation.''

It is useful to check explicitly that the Compton amplitude given in
\eqref{eq:5.74b} obeys the Ward identity.  To do this, replace
$\epsilon_\mu(k)$ by $k_\mu$ or $\epsilon^{*}_\mu(k')$ by $k'_\mu$ and
manipulate the Dirac matrix products.  In either case (after a bit of algebra)
the terms from the two diagrams cancel each other to give zero.

Returning to our derivation of the polarization sum formula \eqref{eq:5.75},
we note that for $k^\mu = (k, 0, 0, k)$, the Ward identity takes the form
\begin{equation}
  k\mathcal{M}^0(k) - k\mathcal{M}^3(k) = 0 .
  \label{eq:5.80}
\end{equation}
Thus $\mathcal{M}^0 = \mathcal{M}^3$, and we have
\begin{equation}
  \begin{aligned}
    \sum_{\epsilon} \epsilon^{*}_\mu(k)\epsilon_\nu(k)\,
      \mathcal{M}^{\mu}(k)\mathcal{M}^{\nu*}(k)
    &= |\mathcal{M}^1|^2 + |\mathcal{M}^2|^2\\
    &= |\mathcal{M}^1|^2 + |\mathcal{M}^2|^2 + |\mathcal{M}^3|^2
      - |\mathcal{M}^0|^2\\
    &= -g_{\mu\nu}\mathcal{M}^{\mu}(k)\mathcal{M}^{\nu*}(k) .
  \end{aligned}
  \label{eq:5.81}
\end{equation}
That is, we may sum over external photon polarizations by replacing
$\sum \epsilon^{*}_\mu\epsilon_\nu$ with $-g_{\mu\nu}$.

Note that this proves (pending our general proof of the Ward identity) that
the unphysical timelike and longitudinal photons can be consistently omitted
from QED calculations, since in any event the squared amplitudes for
producing these states cancel to give zero total probability.  The negative
norm of the timelike photon state, a property that troubled us in the
discussion after Eq.~\eqref{eq:4.132}, plays an essential role in this
cancellation.

\subsection{The Klein-Nishina Formula}

The rest of the computation of the Compton scattering cross section is
straightforward, although it helps to be somewhat organized.  We want to
average the squared amplitude over the initial electron and photon
polarizations, and sum over the final electron and photon polarizations.
Starting with expression \eqref{eq:5.74b} for $\mathcal{M}$, we find
\begin{equation}
  \begin{aligned}
    \frac{1}{4}\sum_{\text{spins}}|\mathcal{M}|^2
    &= \frac{e^4}{4}\ g_{\mu\nu}g_{\rho\sigma}\ \mathrm{tr}\Bigl\{
      (\slashed{p}'+m)\ \Bigl[
        \frac{\gamma^\mu\slashed{k}\gamma^\nu}{2p\cdot k}
        - \frac{\gamma^\nu\slashed{k}'\gamma^\mu}{2p'\cdot k'} \Bigr]\
      (\slashed{p}+m)\ \Bigl[
        \frac{\gamma^\rho\slashed{k}\gamma^\sigma}{2p\cdot k}
        - \frac{\gamma^\sigma\slashed{k}'\gamma^\rho}{2p'\cdot k'} \Bigr]
      \Bigr\} .
  \end{aligned}
  \label{eq:5.82}
\end{equation}
The trace separates into four terms, I, II, III, and IV:
\begin{equation}
  \frac{1}{4}\sum_{\text{spins}}|\mathcal{M}|^2
    = \frac{e^4}{4}\ \Bigl[ \frac{I}{(2p\cdot k)^2}
      + \frac{II + III}{(2p\cdot k)(2p'\cdot k')}
      + \frac{IV}{(2p'\cdot k')^2} \Bigr],
  \label{eq:5.83}
\end{equation}
where I, II, III, and IV are complicated traces.  Note that IV is the same as
I if we replace $k$ with $-k'$.  Also, since we can reverse the order of the
$\gamma$ matrices inside a trace (Eq.~\eqref{eq:5.7}), we see that
$II = III$.  Thus we must work only to compute I and II.

The first of the traces is
\begin{equation}
  I = \mathrm{tr}\bigl[ (\slashed{p}'+m)(\gamma^\mu\slashed{k}\gamma^\nu
    + 2\gamma^\nu p^\nu)(\slashed{p}+m)(\gamma_\mu\slashed{k}\gamma_\nu
      + 2\gamma_\nu p_\nu) \bigr] .
  \label{eq:5.84}
\end{equation}
There are 16 terms inside the trace, but half contain an odd number of
$\gamma$ matrices and therefore vanish.  We must now evaluate the other eight
terms, one at a time.  For example,
\begin{equation}
  \mathrm{tr}[\slashed{p}'\gamma^\mu\slashed{k}\gamma^\nu\slashed{p}
    \gamma_\mu\slashed{k}\gamma_\nu]
  = \mathrm{tr}\bigl[ (-2\slashed{k})\gamma^\nu(-2\slashed{k}')\gamma_\nu
    \bigr]
  = \mathrm{tr}\bigl[ 4\slashed{k}\slashed{k}'(2p\cdot k'-\cdots) \bigr],
\end{equation}
and so on.  Carrying out all the traces and combining terms, we obtain the
standard result for the unpolarized Compton cross section,
\begin{equation}
  \Bigl( \frac{d\sigma}{d\cos\theta} \Bigr)_{\mathrm{lab}}
    = \frac{\pi\alpha^2}{m^2}\ \Bigl(
      \frac{\omega'}{\omega} \Bigr)^2\ \Bigl[
        \frac{\omega'}{\omega} + \frac{\omega}{\omega'}
        - \sin^2\theta \Bigr] ,
  \label{eq:5.94}
\end{equation}
where $\omega$ and $\omega'$ are the initial and final photon energies in the
laboratory frame, related by
\begin{equation}
  \frac{1}{\omega'} - \frac{1}{\omega} = \frac{1}{m}(1-\cos\theta) .
  \label{eq:5.95}
\end{equation}
This is the Klein--Nishina formula.

The singularity of the differential cross section at $\theta\to\pi$ deserves
some attention.  The singular term comes from the square of the
$t$-channel diagram, whose denominator becomes small at backward angles.  Let
us try to understand the physics of this singularity.  The singular term
comes from the square of the $t$-channel diagram,
\begin{equation}
  (e^{-}(p)\ \text{emitting the final photon backward}).
\end{equation}
The amplitude is large at $\theta\approx\pi$ because the denominator of the
propagator is then small ($\sim m^2$) compared to $s$.  To be more precise,
define $\chi$ by
\begin{equation}
  \chi \equiv \pi - \theta .
\end{equation}
We will be interested in values of $\chi$ that are somewhat larger than
$m/\omega$, but still small enough that we can approximate
$1-\cos\chi \approx \chi^2/2$.  For $\chi$ in this range, the denominator is
\begin{equation}
  (p-k')^2 - m^2 = -2p\cdot k' \approx -2\omega^2\Bigl( \frac{m^2}{\omega^2}
    + 1 - \cos\chi \Bigr) \approx -(\omega^2\chi^2 + m^2) .
  \label{eq:5.98}
\end{equation}
This is small compared to $s$ over a wide range of values for $\chi$, hence
the enhancement in the total cross section.

Looking back at Eq.~(5.93), we see that for $\chi$ such that
$m/\omega \ll \chi \ll 1$, the squared amplitude is proportional to
$1/\chi^2$, and hence we expect $\mathcal{M} \propto 1/\chi$.  But we have just
seen that the denominator of $\mathcal{M}$ is proportional to $\chi^2$, so
there must be a compensating factor of $\chi$ in the numerator.  We can
understand the physical origin of that factor by looking at the amplitude for
a particular set of electron and photon polarizations.

Suppose that the initial electron is right-handed.  The dominant term of
Eq.~(5.97) comes from the term that involves
$(\slashed{p}-\slashed{k}')$ in the numerator of the propagator.  Since this
term contains three $\gamma$-matrices in Eq.~(5.97) between the $\bar{u}$
and the $u$, the final electron must also be right-handed.  The amplitude is
therefore
\begin{equation}
  i\mathcal{M} = -ie^2\ \epsilon^{*}_{\mu}(k')\epsilon_\nu(k)\
    \bar{u}_R(p')[\cdots]u_R(p) ,
  \label{eq:5.99}
\end{equation}
where the suppressed factors are the polarization vectors and the numerator of
the propagator.  If the initial photon is left-handed, with
$\epsilon^\mu = (1/\sqrt{2})(0, 1, -i, 0)$, then the combination
$\bar{u}_R(p')\slashed{\epsilon}(k)u_R(p)$ vanishes.  The initial photon must
therefore be right-handed.  Similarly, the amplitude vanishes unless the
final photon is right-handed.  The kinematic situation for this set of
polarizations is shown in Fig.~\ref{fig:5.6}.

\begin{figure}[t]
  \centering
  \includegraphics[width=0.5\textwidth]{fig5.6.png}
  \caption{In the high-energy limit, the final photon is most likely to be
    emitted at backward angles.  Since helicity is conserved, a unit of spin
    angular momentum is converted to orbital angular momentum.}
  \label{fig:5.6}
\end{figure}

Continuing with our calculation, let us consider the numerator of the
propagator in \eqref{eq:5.99}.  For $\chi$ in the range of interest, the
dominant term is
\begin{equation}
  \bm{\sigma}\cdot(\mathbf{p}-\mathbf{k}') \approx \bm{\sigma}\cdot\hat{x}\,
    \omega\chi .
\end{equation}
This is the factor of $\chi$ anticipated above.  It indicates that the final
state is a $p$-wave, as required by angular momentum conservation.
Assembling all the pieces, we obtain
\begin{equation}
  \mathcal{M}(e^{-}_R\gamma_R \to e^{-}_R\gamma_R)
    \approx e^2\sqrt{2E}\ \sqrt{2E'}\ \frac{\omega\chi}{\omega^2\chi^2+m^2}\
    \cdots .
  \label{eq:5.101}
\end{equation}
We would find the same result in the case where all initial and final
particles are left-handed.

Notice that for directly backward scattering, $\chi = 0$, the matrix element
\eqref{eq:5.101} vanishes due to the angular momentum zero in the numerator.
Thus, at angles very close to backward, we should also take into account the
mass term in the numerator of the propagator in Eq.~(5.97).  This term
contains only two gamma matrices and so converts a right-handed electron into
a left-handed electron.  By an analysis similar to the one that led to
Eq.~\eqref{eq:5.101}, we can see that this amplitude is nonvanishing only
when the initial photon is left-handed and the final photon is right-handed.
Following this analysis in more detail, we find
\begin{equation}
  \mathcal{M}(e^{-}_R\gamma_L \to e^{-}_L\gamma_R) = \frac{4e^2 m/\omega}
    {\chi^2 + m^2/\omega^2}
  \label{eq:5.102}
\end{equation}
The reaction with all four helicities reversed gives the same matrix element.

To compare this result to our previous calculations, we should add the
contributions to the cross section from \eqref{eq:5.101} and \eqref{eq:5.102}
and equal contributions for the reactions involving initial left-handed
electrons, and divide by 4 to average over initial spins.  The unpolarized
differential cross section should then be
\begin{equation}
  \begin{aligned}
    \frac{d\sigma}{d\cos\theta} &= \frac{1}{4}\cdot\frac{1}{2E\,2\omega}\,
      \frac{1}{(2\pi)^4}\ \frac{8e^4\chi^2}{(E+\omega)(\chi^2+m^2/\omega^2)^2}
      + \frac{8e^4 m^2/\omega^2}{(\chi^2+m^2/\omega^2)^2}\\
    &= \frac{4\pi\alpha^2}{s}\ \frac{1}{\chi^2+4m^2/s}\ \cdots ,
  \end{aligned}
  \label{eq:5.103}
\end{equation}
which agrees precisely with Eq.~\eqref{eq:5.94}.

The importance of the helicity-flip process \eqref{eq:5.102} just at the
kinematic endpoint has an interesting experimental consequence.  Consider the
process of \emph{inverse Compton scattering}, a high-energy electron beam
colliding with a low-energy photon beam (for example, a laser beam) to
produce a high-energy photon beam.  Let the electrons have energy $E$ and the
laser photons have energy $\omega$, let the energy of the scattered photon be
$E' = yE$, and assume for simplicity that $s = 4E\omega \gg m^2$.  Then the
computation we have just done applies to this situation, with the highest
energy photons resulting from scattering that is precisely backward in the
center-of-mass frame.  By computing $2k\cdot k'$ in the center-of-mass frame
and in the lab frame, it is easy to show that the final photon energy is
related to the center-of-mass scattering angle through
\begin{equation}
  y \approx \frac{1}{1 + (m^2/s)\cdot\text{const}}\ \cdots
    \approx 1 - \frac{m^2}{s}\frac{1}{1-y} .
\end{equation}
Then Eq.~\eqref{eq:5.103} can be rewritten as a formula for the energy
distribution of backscattered photons near the endpoint:
\begin{equation}
  \frac{d\sigma}{dy} = \frac{2\pi\alpha^2}{s}\ \Bigl[
    \frac{1}{(1-y)+m^2/s}\ \cdots\ + \frac{m^2}{s}\frac{1}{(1-y)^2}
      \frac{1}{(1-y)+m^2/s} \Bigr] ,
  \label{eq:5.104}
\end{equation}
where the first term in brackets corresponds to the helicity-conserving
process and the second term to the helicity-flip process.  Thus, for example,
if a right-handed polarized laser beam is scattered from an unpolarized
high-energy electron beam, most of the backscattered photons will be
right-handed but the highest-energy photons will be left-handed.  This effect
can be used experimentally to measure the polarization of an electron beam or
to create high-energy photon sources with adjustable energy distribution and
polarization.

\subsection{Pair Annihilation into Photons}

We can still obtain one more result from the Compton-scattering amplitude.
Consider the annihilation process
\begin{equation}
  e^{+}e^{-} \to 2\gamma ,
\end{equation}
given to lowest order by the diagrams
\begin{equation}
  e^{+}(p_2),\ \gamma(k_1),\ \gamma(k_2) \qquad
  \text{with } p \leftrightarrow \text{Compton substitutions} .
\end{equation}
This process is related to Compton scattering by crossing symmetry; we can
obtain the correct amplitude from the Compton amplitude by making the
replacements
\begin{equation}
  p \to p_1,\qquad p' \to -p_2,\qquad k \to k_1,\qquad k' \to k_2 .
\end{equation}
Making these substitutions in Eq.~(5.87), we find
\begin{equation}
  \frac{1}{4}\sum_{\text{spins}}|\mathcal{M}|^2 = -2e^4\ \Bigl[
    \frac{p_1\cdot k_1}{p_2\cdot k_1} + \frac{p_1\cdot k_2}{p_2\cdot k_2}
      + 2m^2\Bigl( \frac{1}{p_1\cdot k_1} + \frac{1}{p_1\cdot k_2} \Bigr)
      - m^4\Bigl( \frac{1}{p_1\cdot k_1} + \frac{1}{p_1\cdot k_2} \Bigr)^2
    \Bigr] .
  \label{eq:5.105}
\end{equation}
The overall minus sign is the result of the crossing relation \eqref{eq:5.68}
and should be removed.

Now specialize to the center-of-mass frame.  The kinematics is standard, and a
routine calculation yields the differential cross section,
\begin{equation}
  \frac{d\sigma}{d\cos\theta} = \frac{2\pi\alpha^2}{s}\ \Bigl(
    \frac{E}{p} \Bigr)\ \Bigl[
      \frac{E^2 + p^2\cos^2\theta}{m^2 + p^2\sin^2\theta}
      + \frac{2m^2}{m^2+p^2\sin^2\theta}
      - \frac{2m^4}{(m^2+p^2\sin^2\theta)^2}
    \Bigr] .
  \label{eq:5.106}
\end{equation}
In the high-energy limit, this becomes
\begin{equation}
  \frac{d\sigma}{d\cos\theta} \xrightarrow{E\gg m}
    \frac{2\pi\alpha^2}{s}\ \Bigl( \frac{1+\cos^2\theta}{\sin^2\theta}
      \Bigr) .
  \label{eq:5.107}
\end{equation}
The angular dependence of the cross section for $e^{+}e^{-}\to 2\gamma$ at
$E_{\mathrm{cm}} = 29$ GeV, as measured by the HRS collaboration, is shown in
Fig.~\ref{fig:5.7}; the solid line is the lowest-order theoretical prediction,
Eq.~\eqref{eq:5.107}.

\begin{figure}[t]
  \centering
  \includegraphics[width=0.6\textwidth]{fig5.7.png}
  \caption{Angular dependence of the cross section for
    $e^{+}e^{-}\to 2\gamma$ at $E_{\mathrm{cm}} = 29$ GeV, as measured by the
    HRS collaboration, M.~Derrick, et.\ al., \emph{Phys.\ Rev.\ D} \textbf{34},
    3286 (1986).  The solid line is the lowest-order theoretical prediction,
    Eq.~\eqref{eq:5.107}.}
  \label{fig:5.7}
\end{figure}

\problems

\textbf{5.1} Compute the cross section for electron-muon scattering,
$e^{-}\mu^{-}\to e^{-}\mu^{-}$, in the muon rest frame, retaining the electron
mass but sending $m_\mu\to\infty$.

\textbf{5.2} \emph{Bhabha scattering.}  Compute the differential cross section
$d\sigma/d\cos\theta$ for Bhabha scattering, $e^{+}e^{-}\to e^{+}e^{-}$.  You
may work in the limit $E_{\mathrm{cm}} \gg m_e$, in which it is permissible to
ignore the electron mass.  There are two Feynman diagrams; these must be added
in the invariant matrix element before squaring.  Be sure that you have the
correct relative sign between these diagrams.  The intermediate steps are
complicated, but the final result is quite simple.  In particular, you may
find it useful to introduce the Mandelstam variables $s$, $t$, and $u$.  Note
that, if we ignore the electron mass, $s+t+u = 0$.  You should be able to cast
the differential cross section into the form
\begin{equation}
  \frac{d\sigma}{d\cos\theta} = \frac{\pi\alpha^2}{2s}\ \Bigl[ \frac{s^2+u^2}
    {t^2} + \frac{s^2+t^2}{u^2} + \frac{4s^2}{tu} \Bigr] .
\end{equation}
Rewrite this formula in terms of $\cos\theta$ and graph it.  What feature of
the diagrams causes the differential cross section to diverge as $\theta\to 0$?

\textbf{5.3} The spinor product formalism introduced in Problem~3.3 provides
an efficient way to compute tree diagrams involving massless particles.
Recall that in Problem~3.3 we defined spinor products as follows: Let
$u_{L0}$, $u_{R0}$ be the left- and right-handed spinors at some fixed
lightlike momentum $k_0$.  These satisfy
\begin{equation}
  \slashed{k}_0 u_{L0} = \slashed{k}_0 u_{R0} = 0 .
  \label{eq:5.p1}
\end{equation}
(These relations are just the projections onto definite helicity of the more
standard formula $\sum u_0\bar{u}_0 = \slashed{k}_0$.)  Then define spinors for
any other lightlike momentum $p$ by
\begin{equation}
  u_L(p) = \frac{1}{\sqrt{2p\cdot k_0}}\,\slashed{p}\,u_{R0};\qquad
  u_R(p) = \frac{1}{\sqrt{2p\cdot k_0}}\,\slashed{p}\,u_{L0} .
  \label{eq:5.p2}
\end{equation}
We showed that these spinors satisfy $\slashed{p}u_L(p) = \slashed{p}u_R(p)
= 0$; because there is no $m$ around, they can be used as spinors for either
fermions or antifermions.  We defined
\begin{equation}
  s(p_1, p_2) = \bar{u}_R(p_1)u_L(p_2),\qquad
  t(p_1, p_2) = \bar{u}_L(p_1)u_R(p_2),
\end{equation}
and, in a special frame, we proved the properties
\begin{equation}
  t(p_1,p_2) = (s(p_2,p_1))^{*},\qquad
  s(p_1,p_2) = -s(p_2,p_1),\qquad
  |s(p_1,p_2)|^2 = 2p_1\cdot p_2 .
  \label{eq:5.p3}
\end{equation}
Now let us apply these results.
\begin{enumerate}[label=(\alph*)]
  \item To warm up, give another proof of the last relation in Eq.~\eqref{eq:5.p3}
    by using \eqref{eq:5.p1} to rewrite $|s(p_1,p_2)|^2$ as a trace of Dirac
    matrices, and then applying the trace calculus.
  \item Show that, for any string of Dirac matrices,
    \begin{equation}
      \slashed{a}\slashed{b}\cdots = \cdots
    \end{equation}
    where $\mu, \nu, \rho, \ldots = 0, 1, 2, 3$, or 5.  Use this identity to
    show that
    \begin{equation}
      \bar{u}_L(p_1)\gamma^\mu u_L(p_2)
        = \bar{u}_R(p_2)\gamma^\mu u_R(p_1) .
    \end{equation}
  \item Prove the Fierz identity
    \begin{equation}
      \bar{u}_L(p_1)\gamma^\mu u_L(p_2)\,\gamma_\mu = 2\bigl[
        u_L(p_2)\bar{u}_L(p_1) + u_R(p_1)\bar{u}_R(p_2) \bigr],
    \end{equation}
    where $a, b = 1, 2, 3, 4$ are Dirac indices.  This can be done by
    justifying the following statements: The right-hand side of this equation
    is a Dirac matrix; thus, it can be written as a linear combination of the
    16 $\Gamma$ matrices discussed in Section~3.4.  It satisfies
    \begin{equation}
      \gamma^5[M] = -[M]\gamma^5,
    \end{equation}
    thus, it must have the form $V_\mu\gamma^\mu + W_\mu\gamma^\mu\gamma^5$,
    where $V^\mu$ and $W^\mu$ are 4-vectors.  These 4-vectors can be computed
    by trace technology; for example,
    \begin{equation}
      V^\mu = \tfrac{1}{4}\mathrm{tr}[M\gamma^\mu] .
    \end{equation}
  \item Consider the process $e^{+}e^{-}\to\mu^{+}\mu^{-}$, to the leading
    order in $\alpha$, ignoring the masses of both the electron and the muon.
    Consider first the case in which the electron and the final muon are both
    right-handed and the positron and the final antimuon are both left-handed.
    (Use the spinor $v_R$ for the antimuon and $\bar{u}_R$ for the positron.)
    Apply the Fierz identity to show that the amplitude can be evaluated
    directly in terms of spinor products.  Square the amplitude and reproduce
    the result for
    \begin{equation}
      \frac{d\sigma}{d\cos\theta}(e^{-}_R e^{+}_L \to \mu^{-}_R\mu^{+}_L)
    \end{equation}
    given in Eq.~\eqref{eq:5.21}.  Compute the other helicity cross sections
    for this process and show that they also reproduce the results found in
    Section~5.2.
  \item Compute the differential cross section for Bhabha scattering of
    massless electrons, helicity state by helicity state, using the spinor
    product formalism.  The average over initial helicities, summed over
    final helicities, should reproduce the result of Problem~5.2.  In the
    process, you should see how this result arises as the sum of definite-
    helicity contributions.
\end{enumerate}

\textbf{5.4} \emph{Positronium lifetimes.}
\begin{enumerate}[label=(\alph*)]
  \item Compute the amplitude $\mathcal{M}$ for $e^{+}e^{-}$ annihilation into
    2 photons in the extreme nonrelativistic limit (i.e., keep only the term
    proportional to zero powers of the electron and positron 3-momentum).
    Use this result, together with our formalism for fermion-antifermion bound
    states, to compute the rate of annihilation of the $1S$ states of
    positronium into 2 photons.  You should find that the spin-1 states of
    positronium do not annihilate into 2 photons, confirming the symmetry
    argument of Problem~3.8.  For the spin-0 state of positronium, you should
    find a result proportional to the square of the $1S$ wavefunction at the
    origin.  Inserting the value of this wavefunction from nonrelativistic
    quantum mechanics, you should find
    \begin{equation}
      \Gamma = 8.03\times10^9\ \mathrm{sec}^{-1} .
    \end{equation}
    A recent measurement%
    \footnote{D.~W.~Gidley et.\ al., \emph{Phys.\ Rev.\ Lett.} \textbf{49}, 525
    (1982).} gives $\Gamma = 7.994 \pm 0.011\ \mathrm{nsec}^{-1}$; the 0.5\%
    discrepancy is accounted for by radiative corrections.
  \item Computing the decay rates of higher-$\ell$ positronium states is
    somewhat more difficult; in the rest of this problem, we will consider the
    case $\ell = 1$.  First, work out the terms in the $e^{+}e^{-}\to 2\gamma$
    amplitude proportional to one power of the 3-momentum.  (For simplicity,
    work in the center-of-mass frame.)  Since
    \begin{equation}
      \psi(\mathbf{x}=0) \sim \int d^3p\ \tilde{\psi}(\mathbf{p}),
    \end{equation}
    this piece of the amplitude has overlap with $P$-wave bound states.  Show
    that the $S=1$, but not the $S=0$ states, can decay to 2 photons.  Again,
    this is a consequence of $C$.
  \item To compute the decay rates of these $P$-wave states, we need properly
    normalized state vectors.  Denote the three $P$-state wavefunctions by
    \begin{equation}
      \psi_i = \xi_i f(|\mathbf{x}|),\qquad \xi_i = \delta^{ij},
    \end{equation}
    normalized to $\int d^3x\ |\psi_i|^2 = 1$, and their Fourier transforms by
    $\tilde{\psi}_i(\mathbf{p})$.  Show that
    \begin{equation}
      |B(\mathbf{k})\rangle = \sqrt{2M}\int \frac{d^3p}{(2\pi)^3}\
        \tilde{\psi}_i(\mathbf{p})\ \Xi^i\ a^{\dagger}_{\mathbf{k}/2+\mathbf{p}}
        b^{\dagger}_{-\mathbf{k}/2-\mathbf{p}}|0\rangle
    \end{equation}
    is a properly normalized bound-state vector if $\Xi^i$ denotes a set of
    three $2\times2$ matrices normalized to $\mathrm{tr}[\Xi^{i\dagger}\Xi^j]
    = \delta^{ij}$.  To build $S=1$ states, we should take each $\Xi^i$ to
    contain a Pauli sigma matrix.  In general, spin-orbit coupling will split
    the multiplet of $S=1$, $L=1$ states according to the total angular
    momentum $J$.  The states of definite $J$ are given by
    \begin{equation}
      J = 0:\quad \Xi^i = \frac{1}{\sqrt{3}}\sigma^i;\qquad
      J = 1:\quad \Xi^i = -\frac{i}{\sqrt{2}}\epsilon^{ijk}\sigma^k;\qquad
      J = 2:\quad \Xi^i = \frac{1}{\sqrt{2}}h^{ij}\sigma^j,
    \end{equation}
    where $\mathbf{n}$ is a polarization vector satisfying $|\mathbf{n}|^2 = 1$
    and $h$ is a traceless tensor, for which a typical value might be
    $h^{12} = 1$ and all other components zero.
  \item Using the expanded form for the $e^{+}e^{-}\to 2\gamma$ amplitude
    derived in part (b) and the explicit form of the $S=1$, $L=1$, definite-$J$
    positronium states found in part (c), compute, for each $J$, the decay rate
    of the state into two photons.
\end{enumerate}
